% /Users/pikachu/books/payments-book/src/parts/part5/ch28-multicurrency.tex
\chapter{Multi-currency и cross-border}
\label{ch:multicurrency}
\margintoc


\begin{kaobox}[title=Что вы узнаете из этой главы]
  Как принимать архитектурные решения в мультивалютном
  платёжном продукте: кто владеет FX quote, где фиксируется
  settlement currency, как устроены DCC и corridor-based routing,
  чем payout по валютам отличается от payout по мерчантам и почему
  reconciliation в multi-currency почти всегда сложнее, чем сама
  авторизация.
\end{kaobox}

В первых главах книги мы уже разбирали, как работают корреспондентские
счета и межбанковские расчёты, а в части III — как карточные сети
проводят cross-border транзакции. Здесь задача другая: не повторить
теорию rail, а понять, какие решения приходится принимать продукту,
если у него несколько валют, несколько коридоров и несколько мест,
где может возникнуть курсовой риск.

\section{Кто владеет FX quote}
\label{sec:mc-currencies}

Первое архитектурное решение в multi-currency продукте звучит
не как \enquote{какой курс показать}, а как \enquote{кто вообще
имеет право зафиксировать курс}. На практике вариантов четыре.

Первый — курс целиком контролирует эмитент или карточная сеть.
Так живёт стандартный карточный flow: мерчант выставляет цену
в transaction currency, а billing currency определяется только
на стороне эмитента при клиринге. Преимущество — простота
для мерчанта. Недостаток — ни мерчант, ни платформа заранее
не знают финальную сумму списания в валюте карты.

Второй — курс контролирует эквайер или DCC-провайдер.
Тогда держателю карты показывают итоговую сумму в домашней
валюте до авторизации. Это делает цену предсказуемой
для покупателя, но переносит право на margin к стороне,
которая владеет DCC.

Третий — quote принадлежит самой платформе. Так строятся многие
кошельки, travel-платформы и marketplace с собственным FX-слоем:
фронтенд показывает курс, платформа берёт на себя риск его
устаревания на короткое окно, а затем хеджирует или неттит
фактические потоки при settlement.

Четвёртый — quote принадлежит внешнему treasury / FX-провайдеру,
а платёжный продукт лишь привязывает платёж к уже рассчитанной
сделке. Этот вариант типичен для B2B и payout-платформ, где цена
перевода живёт дольше одной checkout-сессии.

\begin{kaobox}[title=На практике]
  Самая опасная архитектура — та, где quote показывается клиенту
  как \enquote{ваш курс}, но фактически продукт не контролирует
  ни источник курса, ни окно его действия. Если владелец FX quote
  не определён явно, спор о потерях на конверсии почти неизбежно
  перейдёт в поддержку и reconciliation.
\end{kaobox}

\section{Где фиксируется settlement currency}
\label{sec:mc-settlement}

После вопроса о quote идёт второй: в какой валюте система
вообще считает свои обязательства. Здесь полезно жёстко развести
три сущности.

\textbf{Transaction currency} — валюта цены, которую видит
покупатель и которую мерчант ожидает на checkout.
\textbf{Billing currency} — валюта, в которой эмитент дебетует
держателя карты. \textbf{Settlement currency} — валюта,
в которой сеть или эквайер фактически рассчитывается с вашей
стороной.

Архитектурная ошибка начинается там, где settlement currency
считается производной от transaction currency по умолчанию.
На практике продукту приходится выбирать одно из двух состояний.

\textbf{Single-currency settlement}. Внутренний реестр обязательств и payouts
ведутся в одной базовой валюте. Это упрощает операционный контур:
один treasury-balance, одна логика комиссий, меньше сценариев
для сверки. Цена этой простоты — постоянная конверсия на входе
или выходе из системы и зависимость от FX markup внешнего провайдера.

\textbf{Multi-currency settlement}. Система хранит обязательства
и расчёты в нескольких валютах. Это сложнее для реестра обязательств и payout-логики,
но позволяет избежать лишней конверсии и не смешивать реальные
валютные позиции в одну синтетическую сумму.

Поэтому хороший вопрос для дизайна звучит не \enquote{поддерживаем ли
мы EUR}, а \enquote{является ли EUR валютой цены, валютой расчёта,
валютой хранения обязательства или всем сразу}. Эти ответы часто
не совпадают.

\section{DCC и cross-border routing}
\label{sec:mc-dcc}

DCC удобнее рассматривать не как маркетинговую опцию на терминале,
а как отдельный routing decision. Как только продукт включает DCC,
он меняет не только UX, но и ownership конверсии: settlement проходит
по одной логике, а billing для клиента — уже по другой.

Это означает три архитектурных последствия.

Во-первых, продукту нужно хранить origin amount, shown amount
и applied FX rate как отдельные поля, а не пытаться вывести одно
из другого постфактум. Без этого невозможно ни объяснить клиенту
списание, ни сверить удержанную DCC-маржу.

Во-вторых, DCC должен быть вписан в routing policy. Для одних
коридоров он экономически оправдан, для других создаёт только
репутационный риск и рост отказов. Внутренняя логика обычно выглядит
так: оценить corridor, BIN, страну эквайера, тип мерчанта, затем
решить, разрешён ли DCC и какой источник quote используется.

В-третьих, DCC нельзя проектировать отдельно от compliance-контуров.
Если opt-in, disclosed markup и audit trail не фиксируются как часть
события, продукт теряет доказательство того, что конверсия была
предложена корректно.

\section{Corridor-based routing и payout architecture}
\label{sec:mc-routing}

Как только у платформы появляется больше одной страны, одного
эквайера или одного платёжного контура выплат, multi-currency превращается
в задачу corridor-based routing: для каждой пары
\enquote{страна покупателя -- страна мерчанта} нужно выбрать
не просто эквайера, а весь маршрут.

Обычно в этом выборе участвуют четыре параметра:
\begin{itemize}
  \item где дешевле авторизация и settlement;
  \item где выше доля одобрений для данного BIN-пула;
  \item в какой валюте мерчант хочет получить payout;
  \item на какой стороне продукт готов держать FX risk.
\end{itemize}
Отсюда вырастают две разные архитектуры payout.

\textbf{Payout by merchant home currency}. Мерчант всегда получает
в домашней валюте, а платформа сама консолидирует и конвертирует
все входящие потоки. Это проще для мерчанта, но сложнее для treasury
платформы.

\textbf{Payout by transaction currency}. Платформа хранит
отдельные валютные карманы и платит мерчанту в той валюте,
в которой пришёл поток. Это уменьшает число конверсий,
но требует, чтобы мерчант умел жить с несколькими валютными
балансами и payout schedule по каждой из них.

\begin{kaobox}[title=На практике]
  Если продукт обещает мерчанту \enquote{один баланс в одной валюте},
  а внутри принимает потоки в пяти валютах, это уже не просто
  платёжная интеграция, а treasury-функция. Её нельзя прятать
  внутри payout-job: нужен явный owner валютной позиции,
  лимитов и источников FX quote.
\end{kaobox}

\section{Reconciliation по валютам}
\label{sec:mc-fx-risk}

Главная operational-боль multi-currency редко возникает
в авторизации. Она возникает позже, когда одна и та же операция
появляется:
\begin{itemize}
  \item в checkout в валюте цены;
  \item в authorizaton log с одной суммой;
  \item в clearing с другой settlement currency;
  \item в payout-отчёте уже после удержания fee и FX spread.
\end{itemize}
Поэтому reconciliation по валютам нужно проектировать как отдельный
слой, а не как частный случай обычной сверки. Практически полезно
держать три правила.

Первое: не терять исходную валютную пару. Внутренний реестр обязательств должен
помнить не только сумму после конверсии, но и origin amount, target amount,
rate source и timestamp фиксации курса.

Второе: считать fee отдельно от FX. Если курсовая разница смешана
с эквайерской комиссией, операционная команда никогда не поймёт,
где именно исчезла маржа.

Третье: сверять payout по валютным карманам, а не по агрегированной
сумме. Сходящаяся общая цифра легко скрывает дыру внутри отдельной
валюты, особенно если одна валюта перекрывает другую по номиналу.

Именно по этой причине multi-currency почти всегда тянет за собой
изменения в реестре обязательств из главы \ref{ch:internal-ledger}: без отдельных
балансов, accrual-модели для fee и связи с payout schedule валютная
сверка становится ручной и хрупкой.

\section{Короткий refresher: correspondent banking и ISO 20022}
\label{sec:mc-correspondent}

Теория correspondent banking подробно разобрана в главе
\ref{ch:banks}, а миграция к \iso{20022} — в главе
\ref{ch:open-banking}. Для архитектурных решений здесь достаточно
удерживать три коротких инварианта.

Первый: cross-border settlement почти всегда проходит по более длинной
цепочке, чем авторизация, и поэтому живёт в другом тайминге.
Approval за секунду не означает payout в тот же день.

Второй: расчётный банк, карточная сеть и FX-провайдер могут быть
тремя разными сторонами. Значит, продукту нужен явный mapping,
кто является owner каждой стадии, а не абстрактное
\enquote{деньги идут через SWIFT}.

Третий: переход отрасли на \iso{20022} делает сообщения богаче,
но не убирает продуктовые решения. Structured fields облегчают
screening и reconciliation, однако сами по себе не решают,
кто несёт FX risk и в какой валюте должны жить внутренние обязательства.

\begin{chaptersummary}
\begin{itemize}
  \item Multi-currency начинается с ownership FX quote: курс может
        контролировать эмитент, эквайер/DCC, сама платформа
        или внешний treasury-провайдер.
  \item Transaction currency, billing currency и settlement currency
        нельзя смешивать: продукт должен явно понимать, какая валюта
        относится к цене, какая — к списанию, какая — к расчёту.
  \item DCC — это routing decision с собственным compliance-аудитом,
        а не просто красивое отображение суммы в домашней валюте.
  \item Corridor-based routing объединяет долю одобрений, стоимость,
        payout-валюту и ownership FX risk в одно архитектурное решение.
  \item Основная сложность multi-currency живёт в реестре обязательств и сверке:
        origin amount, converted amount, source of rate и FX fee
        должны храниться отдельно.
  \item Correspondent banking и \iso{20022} задают инфраструктурный
        контур cross-border, но не снимают продуктового выбора
        о том, где держать валютный риск и в какой валюте считать
        обязательства.
\end{itemize}
\end{chaptersummary}
